\chapter{Primitivos}

\section{Origen canónico}
\begin{definition}
El símbolo \subz\ denota el origen canónico del sistema. Toda coordenada, medida u operación se interpreta de manera relativa a \subz.
\end{definition}

\begin{observation}
La introducción de \subz\ fija una referencia interna de arranque estructural. No valida por sí sola relaciones posteriores que dependan de desarrollos adicionales.
\end{observation}

\section{Resolución mínima}
\begin{definition}
El símbolo \eps\ denota la resolución mínima legible del sistema, esto es, el primer paso positivo distinguible.
\end{definition}

\begin{observation}
La función de \eps\ es impedir el colapso entre ausencia y primer desplazamiento legible. Su introducción fija granularidad mínima, pero no implica por sí sola estabilidad global de regímenes posteriores.
\end{observation}

\section{Ventana local}
\begin{definition}
La ventana \Win\ declara el dominio local de lectura dentro del cual una configuración puede ser distinguida, comparada y evaluada sin ambigüedad inmediata.
\end{definition}

\begin{observation}
La noción de ventana es local: organiza lectura, comparación y contraste dentro de un dominio declarado, sin agotar por sí sola la arquitectura completa del sistema.
\end{observation}

